\section{ARIMA：非平稳序列的差分处理}
\label{sec:06-Random-Processes/06-ARMA-and-ARIMA/03}

你拿到一支股票的多年日收盘价：从几十元缓慢攀升到几百元。用 AR(1) 拟合，估计出 \(\hat\phi\approx 0.998\)，几乎就是 1。你试着画 ACF——它衰减得像是永远到不了零。残差里还残留着缓慢的漂移结构。

这不是 ARMA 本身的问题。问题的根源在数据：序列的均值随时间持续漂移，方差也可能逐年扩大。ARMA 的平稳性假设被彻底违反了。

\subsection{随机游走：当记忆永不消退}

AR(1) 模型 \(X_t=\phi X_{t-1}+\varepsilon_t\) 在 \(\phi=1\) 处发生了质的改变。\(\phi<1\) 时，扰动的影响随时间指数衰减——昨天的一个冲击，对一个月后的影响已经微乎其微。\(\phi=1\) 时：

\[
X_t=X_{t-1}+\varepsilon_t,
\qquad
X_t=X_0+\sum_{j=1}^{t}\varepsilon_j.
\]

每一个 \(\varepsilon_j\) 的影响都永不衰减地被累积进后续的所有观测。展开式中，\(X_t\) 是初始值加上所有历史创新的和——这个和随着 \(t\) 增大，方差线性增长：

\[
\Var(X_t)=t\sigma^2.
\]

它不像平稳过程那样有一个可以回归的均值。最好的预测是``明天等于今天''，因为任何偏离都是零均值噪声，且没有均值回复力。ACF 也表现得不像平稳过程：

\[
\rho_k=\sqrt{1-\frac{k}{t}}\approx 1\quad(\text{对 }k\ll t).
\]

这就是\textbf{单位根过程}：AR 特征多项式 \(1-\phi z=0\) 的根恰好落在单位圆上（\(z=1\)）。它标志着平稳与非平稳的分界线。在 \(\phi\) 穿越 1 的过程中，过程的定性行为发生了一次相变。

单位根检验（ADF、KPSS）试图从有限样本中判断是否存在单位根。ADF 的原假设是``序列有单位根''，拒绝原假设后才支持平稳备择。KPSS 的原假设反过来是``序列平稳''，拒绝原假设意味着数据不支持平稳假设。两种检验的逻辑互补——ADF 告诉你``数据是否不像有单位根''，KPSS 告诉你``数据是否不像平稳''。两者的结论有时会因为样本长度不足、存在结构突变或趋势项的设定差异而出现分歧。机械地看``哪个检验通过就用哪个结果''是有风险的：一个有突变点的平稳序列可能被 ADF 错判为单位根，一个差分平稳序列在短样本里可能被 KPSS 接受为平稳。两个检验结合使用，分歧本身就是进一步诊断的信号。

\subsection{差分让问题回到平稳框架}

面对非平稳序列，有两条截然不同的处理路线。

\textbf{差分平稳（difference stationary，DS）：} 序列有单位根，但差分后平稳。定义一阶差分

\[
\Delta X_t=X_t-X_{t-1}.
\]

若 \(X_t\) 是随机游走，\(\Delta X_t=\varepsilon_t\) 已经是白噪声——一阶差分直接把单位根消掉了。如果一阶差分后仍不平稳，继续差分：

\[
\Delta^2 X_t=\Delta(\Delta X_t)=X_t-2X_{t-1}+X_{t-2}.
\]

\textbf{趋势平稳（trend stationary，TS）：} 序列没有单位根，但有一个确定性的时间趋势。模型为

\[
X_t=\alpha+\beta t+Y_t,
\]

其中 \(Y_t\) 是零均值的平稳过程。这种情况下，正确的处理不是差分，而是先估计并移除趋势，再对残差 \(Y_t\) 拟合 ARMA。

混淆 DS 与 TS 的代价很高：

\begin{itemize}
\tightlist
\item
  对 DS 序列做去趋势（拟合一条时间直线然后取残差）：单位根仍然残留在残差中。残差的 ACF 仍极慢衰减，在此基础上做的回归会出现虚假的统计显著——一个完全随机的单位根序列也可能``显著地''与另一个无关的单位根序列``相关''。
\item
  对 TS 序列做差分：差分后的序列会引入一个不可逆的 MA 单位根，模型变成过度差分。TS 序列的平稳残差被差分一次后，变成了一个非平稳的 MA(1) 过程，ACF 会出现特征性的 \(-0.5\) 等等模式。
\end{itemize}

区分 DS 与 TS 不能仅凭肉眼。DS 序列的方差随时间增长，TS 序列在去趋势后方差恒定。DS 序列的 ACF 衰减极慢（几十阶后仍显著为正），TS 序列去趋势后的 ACF 正常衰减。在实践中，金融价格的对数序列通常被判为 DS——对数收益率（差分后的序列）表现出平稳特征——而某些宏观变量可能更接近 TS。

\subsection{ARIMA(p,d,q)：把差分和 ARMA 统一在一个符号下}

一旦确定了差分阶数 \(d\)，后续就是标准的 ARMA 建模。ARIMA(p,d,q) 用算子形式统一表达：

\[
\phi(B)(1-B)^d X_t=\theta(B)\varepsilon_t,
\]

其中 \((1-B)^d\) 是差分算子。\(B\) 是后移算子：\(BX_t=X_{t-1}\)。\((1-B)^d\) 展开后是 \(d\) 阶差分。\(\phi(B)\) 和 \(\theta(B)\) 是 AR 和 MA 的特征多项式，它们的根现在要求落在单位圆外——但这是对差分后的序列而言，不是对原始序列。

几个常见的特例承载了 ARIMA 框架的大部分应用场景：

\begin{itemize}
\tightlist
\item
  ARIMA(0,1,0)：随机游走——最简单的非平稳模型，也是许多金融序列的基线。
\item
  ARIMA(1,1,0)：差分后的 AR(1)——股价收益率常有微弱的自相关。
\item
  ARIMA(0,1,1)：简单指数平滑的 ARIMA 等价形式。指数平滑预测等价于一个特定参数约束下的 ARIMA(0,1,1) 模型。
\item
  ARIMA(0,2,2)：Holt 线性趋势方法的 ARIMA 对应。二阶差分配合 MA(2) 捕捉了局部线性趋势加短期波动。
\end{itemize}

差分阶数 \(d\) 的选择原则上是自下而上的：从 \(d=0\) 开始，检查残差 ACF 的衰减速度和单位根检验结果。如果 ACF 衰减极慢或单位根检验无法拒绝，增加 \(d\) 并重复。实践中 \(d\le 2\) 几乎覆盖了所有真实序列——\(d=2\) 已经对应线性趋势的随机游走（趋势的斜率本身是随机游走），\(d>2\) 罕见且常常意味着模型中遗漏了结构突变或确定性趋势。

\subsection{Box-Jenkins：建模是迭代验证，不是一次性拟合}

Box-Jenkins 方法论的贡献，不在于发明了 ARIMA，而在于把时间序列建模从``选一个模型，估计参数，结束''变成了一次有反馈闭环的探索流程。四步循环：

\begin{enumerate}
\tightlist
\item
  \textbf{识别（identification）}：检查平稳性，确定 \(d\)；用样本 ACF 和 PACF 的模式初步判断 \(p\) 和 \(q\) 的范围——AR(\(p\)) 的 ACF 拖尾而 PACF 在 \(p\) 阶截尾，MA(\(q\)) 则相反；用 AIC 或 BIC 在几个候选模型中做初步筛选。
\item
  \textbf{估计（estimation）}：对选定的 \((p,d,q)\) 估计参数。纯 AR 可用 Yule-Walker 方程，ARMA/ARIMA 通常用最大似然或条件最小二乘。MLE 需要处理初始观测的分布，短序列时初始化方式会影响估计。
\item
  \textbf{诊断（diagnosis）}：检查残差是否是白噪声——Ljung-Box 检验若干个滞后上的联合自相关性是否接近零。检查残差的正态性和同方差性（QQ 图、残差-时间图）。如果残差仍有结构（显著的季节模式、方差聚集、异常值），回到识别步。
\item
  \textbf{预测（forecasting）}：用估计好的模型生成未来值及其预测区间。ARIMA 的预测是最小均方误差意义下的最优预测，但预测区间依赖于创新正态假设——肥尾的创新会使实际覆盖概率低于名义值。
\end{enumerate}

诊断步是这个循环的心脏。如果诊断不通过就回到识别——这个过程可能重复两三轮——而不是在第一次参数估计后就报告结果。Box-Jenkins 不是一套公式，是一套工作纪律。

\subsection{什么时候 ARIMA 不够用}

ARIMA 有效的前提是线性关系、差分足以获得平稳、创新是独立同分布的。当这些前提之一不成立时，有对应的扩展方向：

\begin{itemize}
\tightlist
\item
  \textbf{结构突变}：参数或数据生成机制随时间改变。ARIMA 会把突变误认为非平稳，过度差分。需要变点检测或状态空间模型（允许参数缓慢演化）。
\item
  \textbf{长记忆}：ACF 衰减不是指数级而是幂律级，差分阶数 \(d=1\) 不够，但 \(d=2\) 又过度。ARFIMA 允许分数阶差分 \(d\in(0,0.5)\)，让记忆长度在两个整数阶之间连续调节。
\item
  \textbf{非线性依赖}：条件均值的演变不能用线性的 AR 结构捕捉。TAR（门限自回归）允许不同区间有不同的 AR 系数；STAR 让系数平滑过渡。
\item
  \textbf{异方差}：创新方差本身有序列依赖——大波动倾向于聚集出现。GARCH 族模型直接对方差建模。
\item
  \textbf{多变量}：多个序列之间存在动态交互和共同趋势。VAR（向量自回归）是多变量 AR 的直接推广。如果多个序列各自有单位根但它们的某个线性组合是平稳的，就涉及协整——长期均衡关系与短期修正动态同时存在。
\end{itemize}

ARIMA 的强处不在于它能处理一切序列，而在于它的假设足够明确：你可以系统地检验每一个前提，并且当检验不通过时，你能指出应该朝哪个方向扩展、需要额外引入什么结构。它不是序列的近义词，是序列分析的一条基线。

下一篇把 ARIMA 放在更大的视野中——与状态空间模型和机器学习方法并列——看它们在什么样的数据条件下各自自然地成为首选。
